

Across software domains, from general-purpose systems to long-lived industrial control codebases, a vast body of legacy code lacks formal specifications, i.e., precise, verifiable descriptions of program behavior such as preconditions, postconditions, and loop invariants. Automating specification generation is therefore essential for program comprehension, reuse, and verification. The task is challenging due to the interplay of complex control structures (nested branches, loops, function calls) and data representations (structs, pointers, inductive datatypes). In this work, we tackle \emph{goal-independent} functional specification generation: inferring verifier-checked preconditions, postconditions, and loop invariants that capture functional behavior without requiring a user-supplied assertion as the synthesis target. The goal is not to prove one preselected property, but to recover verifier-checked specifications that capture a program's functional behavior, whether or not symbolic execution can derive them exactly.

\DIFaddbegin
Existing work on program specification generation has largely focused on loop invariants. Traditional symbolic methods such as constraint solving~\cite{DBLP:conf/cav/ColonSS03,DBLP:journals/pacmpl/LiuFYSL22}, Craig interpolation~\cite{FiB_Squeezing_loop,DBLP:conf/cade/GanDXZKC16,Nonlinear_Craig_Interpolant_Generation}, recurrence analysis~\cite{beyer2024decomposing,wonisch2012predicate}, shape analysis~\cite{calcagno2009compositional}, and dynamic methods such as Daikon~\cite{Daikon} are each restricted to particular invariant forms or program classes; recent ML-augmented or SMT-augmented methods~\cite{Code2Inv_a_Deep_Learning_Framework_for_Program_Verification,Loop_Invariant_Inference_through_SMT_Solving_Enhanced_Reinforcement_Learning,DBLP:journals/pacmse/CaoWXYWCM25} are similarly confined to numeric programs. LLM-based approaches have begun to broaden the program coverage, but tend to produce specifications that are either ill-formed, too weak, or both, and recent integrations with static/formal techniques~\cite{autospec,Enhancing_Automated_Loop_Invariant_Generation_for_Complex_Programs_with_Large_Language_Models,SpecGen,Wang2025ATO} only partially address this. \autospec~\cite{autospec} and Preguss~\cite{Wang2025ATO} are goal-dependent (requiring a predefined verification target or RTE target) and degrade on complex data structures; \specgen~\cite{SpecGen} is goal-independent but only in the Java/JML setting, and ``verifier-accepted'' alone says nothing about strength. These limitations leave open the problem of generating goal-independent, semantically informative, verifier-checked specifications for C programs with loops, calls, pointers, structs, and inductive data.
\DIFaddend


\DIFaddbegin
To address this open problem, we propose \toolname, a framework that integrates symbolic execution, LLM-based synthesis, and formal verification. The key idea is to let symbolic execution drive the process. The analysis proceeds through the program and pauses at program boundaries (loops, calls, and recursive constructs), turning one whole-function specification task into a sequence of smaller local synthesis problems. On loop-free segments fully handled by the flattened-memory symbolic executor, symbolic execution computes an exact path-sensitive postcondition, which is the strongest postcondition within that model; outside this envelope, symbolic execution provides partial semantic anchors for LLM synthesis. At each boundary, the symbolic state, i.e. the symbolic store of variable values, the path conditions, and the types of the in-scope variables, collectively the \emph{local symbolic context}, is injected into invariant templates and LLM prompts. When a candidate fails, formal verification feedback is combined with symbolic information to refine the failed candidate rather than re-sampling blindly.
\DIFaddend

\DIFaddbegin
We implement \toolname\ on top of QCP's symbolic execution engine~\cite{vst-a,wu2025qcppracticalseparationlogicbased} and Frama-C/WP formal verifier~\cite{cuoq2012frama}.
We perform a comprehensive evaluation of \toolname, including an ablation study that confirms the contribution of each component. We evaluate on benchmark families from prior invariant- and specification-generation work, covering numerical loops, multi-/nested loops, pointers, structs, and singly linked lists. On numerical benchmarks, \toolname\ approaches the performance of tools specialized for numerical programs and consistently outperforms recent LLM-based approaches such as \autospec~\cite{autospec} and ACInv~\cite{Enhancing_Automated_Loop_Invariant_Generation_for_Complex_Programs_with_Large_Language_Models} (a static-analysis-guided LLM invariant generator). On the pointer-, struct-, and linked-list-heavy families that prior LLM-based tools handle poorly, it outperforms \autospec\ by a clear margin. We further construct \textbf{\toolname-400}, a 400-program benchmark with all assertions removed, to stress-test the goal-independent setting at scale, and use it to analyze the validity, strength, failure modes, and cost of generated specifications; here \toolname\ substantially outperforms the single-shot LLM and \autospec\ baselines, and a pairwise judge ranks its specifications as stronger. Together these results show that the symbolic-execution-guided design generalizes from the goal-dependent to the goal-independent setting without specialization.
\DIFaddend

\DIFaddbegin
In summary, the key contributions of this work are:
\begin{itemize}
\item \textbf{Symbolic-execution-guided specification synthesis}: 
We present a framework that uses path-sensitive symbolic states to guide LLM-based specification generation. On segments fully covered by symbolic execution, these states provide exact semantic summaries; outside the fully covered fragment, the same symbolic state still supplies local symbolic context to anchor templates and prompts for LLM synthesis.
\item \textbf{Verifier-guided iterative refinement}: We design a multi-layer refinement mechanism that classifies Frama-C/WP errors into diagnostic categories and applies targeted repair strategies, progressively guiding the LLM toward verifier-checked specifications instead of re-sampling blindly.
\item \textbf{Goal-independent evaluation}: We construct \textbf{\toolname-400}, a benchmark of 400 C programs with assertions removed, and use it to conduct what we believe is the first empirical study of goal-independent specification generation, evaluating validity, strength, failure modes, and per‑program specification generation cost.
\end{itemize}
\DIFaddend


\DIFaddbegin
The rest of the paper is organized as follows.
Section~\ref{sec:motivation} introduces a motivating example and gives an overview of our approach.
Section~\ref{sec:summarygeneration} and Section~\ref{sec:loopinvariant} present the core techniques for pre-/postcondition specification generation and loop-invariant generation, respectively.
Section~\ref{sec:refinement} describes the verifier-guided refinement procedure, and Section~\ref{sec:impl} presents the system design and implementation.
Section~\ref{sec:experiments} reports the empirical evaluation under established benchmark protocols.
Section~\ref{sec:large-scale-experiment} presents a goal-independent study on the \textbf{\toolname-400} benchmark.
Section~\ref{sec:case} gives representative case studies.
Section~\ref{sec:limitations} discusses limitations and open challenges, and Section~\ref{sec:threats} discusses threats to validity.
Section~\ref{sec:relatedwork} surveys related work, and Section~\ref{sec:conclusion} concludes.
Finally, the data-availability statement provides links to the datasets, implementation code, LLM configurations, and supplementary results.
\DIFaddend
